\section{Тест производительности}

\subsection{Методика}

Для каждого размера $N$ генерируется набор уникальных ключей из строчных букв латинского
алфавита случайной длины от 1 до 16 символов. Измеряются три операции отдельно:
вставка $N$ ключей, поиск всех $N$ ключей, удаление всех $N$ ключей. Сравнение
проводится с \texttt{std::map}~--- сбалансированным деревом поиска из STL с
гарантированной сложностью $O(\log n)$ на все операции.

Время каждой операции замеряется с помощью \texttt{std::chrono::high\_resolution\_clock}.
Генерация данных и ввод-вывод в замер не входят.

\subsection{Результаты}

\begin{table}[!ht]
\centering
\begin{tabular}{|r|r|r|r|r|r|r|}
\hline
\rowcolor{lightgray}
& \multicolumn{3}{c|}{PATRICIA, мкс} & \multicolumn{3}{c|}{\texttt{std::map}, мкс} \\
\hline
\rowcolor{lightgray}
$N$ & Insert & Find & Remove & Insert & Find & Remove \\
\hline
1\,000    & 348    & 208    & 465    & 214    & 161    & 303    \\
\hline
10\,000   & 4\,731 & 3\,246 & 8\,264 & 4\,251 & 2\,405 & 3\,711 \\
\hline
100\,000  & 86\,469 & 75\,838 & 100\,506 & 87\,293 & 73\,691 & 81\,991 \\
\hline
500\,000  & 731\,836 & 695\,895 & 764\,450 & 766\,442 & 748\,168 & 785\,446 \\
\hline
1\,000\,000 & 1\,842\,835 & 2\,022\,130 & 2\,038\,648 & 1\,953\,836 & 1\,985\,049 & 2\,322\,171 \\
\hline
\end{tabular}
\caption{Время операций в микросекундах (замер в Docker, gcc 14, -O2)}
\end{table}

\begin{alltt}
# python3 lab2/benchmark/generator.py 1000
# ./build/lab2/lab2_benchmark < tests.txt
N=1000
PATRICIA  Insert=348us Find=208us Remove=465us
std::map  Insert=214us Find=161us Remove=303us
# python3 lab2/benchmark/generator.py 10000
# ./build/lab2/lab2_benchmark < tests.txt
N=10000
PATRICIA  Insert=4731us Find=3246us Remove=8264us
std::map  Insert=4251us Find=2405us Remove=3711us
# python3 lab2/benchmark/generator.py 100000
# ./build/lab2/lab2_benchmark < tests.txt
N=100000
PATRICIA  Insert=86469us Find=75838us Remove=100506us
std::map  Insert=87293us Find=73691us Remove=81991us
# python3 lab2/benchmark/generator.py 500000
# ./build/lab2/lab2_benchmark < tests.txt
N=500000
PATRICIA  Insert=731836us Find=695895us Remove=764450us
std::map  Insert=766442us Find=748168us Remove=785446us
# python3 lab2/benchmark/generator.py 1000000
# ./build/lab2/lab2_benchmark < tests.txt
N=1000000
PATRICIA  Insert=1842835us Find=2022130us Remove=2038648us
std::map  Insert=1953836us Find=1985049us Remove=2322171us
\end{alltt}

\subsection{Анализ результатов}

На малых объёмах ($N \leq 10\,000$) \texttt{std::map} обычно быстрее: стандартная
реализация хорошо оптимизирована, а накладные расходы PATRICIA на выделение узлов,
побитовые проверки и хранение строк ещё не окупаются.

На среднем размере ($N = 100\,000$) результаты близки: вставка и поиск отличаются
незначительно, удаление у PATRICIA остаётся дороже из-за более сложного восстановления
ссылок после удаления узла.

На больших размерах ($N \geq 500\,000$) PATRICIA начинает выигрывать на части операций:
при $N = 500\,000$ поиск и удаление быстрее \texttt{std::map}, а при $N = 10^6$
PATRICIA быстрее по вставке и удалению. Это соответствует идее сжатого trie: путь
задаётся критическими битами ключей, а не полным лексикографическим сравнением строк
на каждом уровне сбалансированного дерева.

\pagebreak
